Het was weer een zondag.
Zondag, de dag van macht en zelfreflectie, benadrukt Svenska’s dominantie en ambitie.

vanuit mijn perspectief draait alles om mij

Monthy Python ridder je komt er niet langs is Mara

Svenska met oogkleppen
hield erg van opscheppen.

Men noemde dat ‘snakken’,
want erin gebakken,

zat snakken naar aandacht,
dat haar bestaansrecht bracht.

Ze prees zichzelf aldoor.
Getroost door het gehoor

naar haar doldwaze trots.
Vanaf haar aperots

veinsde ze schranderheid
Doch haar onwetenheid

vermomd als een vals pad
naar het warme breinbad

van hogere kennis
kwam als heiligschennis.

"Foolish pride" (dwaze trots) is een concept dat in zowel de psychologie als de filosofie wordt onderzocht als een destructieve vorm van eigenwaan, vaak voortkomend uit onwetendheid . Het wordt gezien als een "vals pad" naar spiritualiteit of hogere kennis, waarbij de persoon te trots is om de eigen beperkingen of onwetendheid te erkennen.
Hier zijn de belangrijkste perspectieven:
Filosofisch perspectief: In de filosofie wordt "foolish pride" vaak gekoppeld aan een gebrek aan phronēsis (praktische wijsheid) of zelfkennis, vergelijkbaar met de Socratische notie dat ware wijsheid begint met het erkennen van de eigen onwetendheid. Het kan worden gezien als een vorm van "hoogmoed" die de mens blind maakt voor de werkelijkheid.
Psychologisch perspectief: In de psychologie wordt dit vaak besproken in de context van zelfrealisatie en de mythen die mensen creëren om hun ego te beschermen. Het kan leiden tot een onvermogen om fouten toe te geven, wat de persoonlijke ontwikkeling en "zorg voor het zelf" (zoals besproken door Foucault) belemmert.
Spiritualiteit \& Ethiek: Het wordt beschouwd als een obstakel op het pad naar ware kennis of spirituele groei, omdat de trots het individu verhindert om nederigheid te tonen en open te staan voor nieuwe inzichten.
Literatuurverwijzing: Een specifiek boek over dit onderwerp is "Foolish Pride" van Chad E. Peterson.

dwaze trots.

Iemand die zichzelf overdreven prijst, troost of emotioneel oppept, vaak op een manier die narcisme of zelfgenoegzaamheid suggereert.

Ze gebruikte snakken
omdat ze snakte

het martelaarsyndroom. mensen die zo'n houding aannemen, proberen via het aannemen van een slachtofferrol aandacht en genegenheid te verkrijgen.

het martelaarsyndroom en narcisme zijn beide persoonlijkheidsuitdrukkingen die te maken hebben met de behoefte aan aandacht, maar ze verschillen in hun kernmotivaties. terwijl narcisme vaak gepaard gaat met een overdreven gevoel van eigenwaarde en een verlangen naar bewondering, draait het martelaarsyndroom om het innemen van een slachtofferrol om sympathie en steun te verkrijgen. soms kunnen de twee overlappen, maar ze spreken in principe over verschillende mechanismen van zelfpresentatie.

er is geen algemeen erkende term voor de overlap, maar soms wordt dit gedrag omschreven als covert narcisme, vulnerable narcisme of narcistisch slachtofferschap. in deze context presenteert iemand zich als slachtoffer om op een subtiele, vaak introverte manier narcistische behoeften aan erkenning en bewondering te vervullen.

iemand met deze eigenschappen herken je vaak aan een patroon waarin ze voortdurend de slachtofferrol aannemen, maar tegelijkertijd een stille behoefte hebben aan bewondering en bevestiging. ze lijken vaak bescheiden en introvert, maar praten herhaaldelijk over hun persoonlijke lijden en belemmeringen. hierbij een aantal herkenningspunten:

ze benadrukken steevast hoe hard het leven hen beproeft en hoezeer ze onrecht ervaren, ook als de omstandigheden minder zwaar lijken.
ze geven zelden verantwoordelijkheid voor hun eigen gedrag en wijzen vaak de schuld aan externe factoren of andere mensen.
ze lijken op het eerste gezicht kwetsbaar en zelfkritisch, maar verwachten in feite medelijden en lof van hun omgeving.
ze reageren vaak gevoelig of passief-agressief op kritiek, omdat dit de bevestiging ondermijnt die ze zo hard zoeken.
ze vertellen verhalen over hun lijden die vaak een zekere dramaturgie bevatten, wat kan wijzen op een onderliggende behoefte aan aandacht.
ondanks hun slachtofferrol ontbreekt vaak de oprechte empathie voor anderen; gesprekken draaien grotendeels om hun eigen ervaringen.
dit patroon, soms aangeduid als covert of vulnerable narcisme, maakt het lastiger om de persoon in één oogopslag te doorgronden, omdat hun uitstraling minder overdreven is dan bij openlijk narcistisch gedrag. ze balanceren tussen het zoeken naar medelijden en het verlangen naar bewondering, waardoor hun gedrag subtiel en genuanceerd kan overkomen.

dit gedrag kenmerkt zich doordat de persoon een theatrale combinatie hanteert van slachtofferschap en een verlangen naar publieke bewondering. de persoon presenteert zichzelf als onoverwinnelijk in het licht van tegenslagen, maar doet dit niet in een intieme setting waar de nabijheid en authenticiteit centraal staan, maar juist met de bedoeling een zo breed mogelijk publiek te imponeren. dit kan duiden op een vorm van covert of vulnerable narcisme, waarbij de nadruk op eigen ‘dapperheid’ en het vermogen om tegenspoed te overwinnen een middel is om erkenning en bewondering te verkrijgen. de publieke presentatie dient dan minder om werkelijke kwetsbaarheid te delen en meer als een podium waarop men zichzelf als bijzonder en bewonderenswaardig kan etaleren, wat een ingewikkelde dynamiek creëert tussen het aannemen van een slachtofferrol en het zoeken naar narcistische bevestiging.

Solipsisme is de filosofische en psychologische theorie dat alleen het eigen bewustzijn met zekerheid bestaat. Alles buiten het eigen 'ik' – andere mensen, de fysieke wereld, objecten – wordt beschouwd als een creatie of projectie van de eigen geest. De term komt van het Latijnse solus (alleen) en ipse (zelf).

Wikikids +2
Hieronder volgt een gedetailleerd overzicht van solipsisme in filosofie en psychologie:
Solipsisme in de Filosofie
Fundament (Descartes): Het is geworteld in radicaal scepticisme en het Cartesianisme, samengevat in Cogito, ergo sum ("Ik denk, dus ik ben"). De enige zekerheid is het bestaan van de denkende geest zelf.
Het Probleem van Andere Geesten: Solipsisme stelt dat men nooit direct kan weten of anderen een bewustzijn hebben. Anderen kunnen gezien worden als 'automaten' of projecties.
Variëteiten: Er wordt onderscheid gemaakt tussen metafysisch solipsisme (enkel de eigen geest bestaat echt), epistemologisch solipsisme (men kan alleen de eigen geest kennen) en methodologisch solipsisme (gebruikt als sceptische methode).
Status: Hoewel invloedrijk in sceptische debatten, is er nauwelijks een filosoof die het solipsisme als een serieuze, letterlijke waarheid verdedigt.

Filosofie Magazine +4

Solipsisme in de Psychologie en Ontwikkeling
Ontwikkelingspsychologie (Piaget): Jean Piaget suggereerde dat pasgeborenen zich in een staat van 'oorspronkelijk solipsisme' bevinden. Zij maken initieel geen onderscheid tussen zichzelf en de buitenwereld.
Psychopathologie: Solipsistische gedachten kunnen voorkomen tijdens psychoses, waarbij de grens tussen de interne belevingswereld en de externe realiteit vervaagt.
Angststoornis: Het kan leiden tot existentiële angst of derealisatie, waarbij de betrokkene het gevoel heeft dat niets echt is.

PsychoseNet +2

Verschil met Andere Begrippen
Solipsisme vs. Narcisme: Waar narcisme een obsessie is met het eigen 'zelf' (superioriteit), is solipsisme de (theoretische) weigering of onmogelijkheid om te erkennen dat er iets anders dan het zelf bestaat.

Reddit

Kritiek
Solipsisme is in de praktijk onhoudbaar: de wereld toont zich en is niet slechts een eigen creatie.
Het wordt vaak gebruikt als een gedachte-experiment (zoals in "leef ik in een droom of simulatie?") in plaats van een levenswijze.

www.faidros.nl +2

"Rock bottom" (het absolute dieptepunt) is een complex concept dat psychologische en filosofische inzichten combineert. Het wordt vaak gezien als een noodzakelijk keerpunt.

Medium
Hier zijn de belangrijkste aspecten:
Psychologische inzichten
Keerpunt en Wake-up Call: "Rock bottom" is het punt waarop de huidige levensstijl onhoudbaar is geworden, wat kan leiden tot een noodzakelijke verandering of een "wake-up call".
Emotionele verwerking: Het proces kan pijnlijk zijn (denk aan relatiebreuken of persoonlijk faillissement), maar biedt tegelijkertijd de gelegenheid om te leren en emotioneel te groeien.
Persoonlijke betekenis: Wat voor de één "rock bottom" is, hoeft dat voor de ander niet te zijn; het is een subjectieve ervaring van falen of verlies.
Wederopbouw: Hitting rock bottom kan de basis vormen om het leven opnieuw op te bouwen.

Medium +1

Filosofische perspectieven
Existentialisme: Het dieptepunt dwingt een individu tot zelfreflectie en het nemen van verantwoordelijkheid voor de eigen keuzes en toekomst.
Positieve desintegratie: De ervaring van een crisis of dieptepunt kan leiden tot een hogere vorm van persoonlijke ontwikkeling (positieve desintegratie).
Kritiek op het concept: Er zijn filosofische kanttekeningen bij de noodzaak van "rock bottom", waarbij wordt betwijfeld of een dergelijk extreem dieptepunt noodzakelijk is voor verandering.
Betekenisgeving: Filosofisch gezien kan het dieptepunt worden gezien als een moment om opnieuw te definiëren wat belangrijk is in het leven.

www.positievedesintegratie.nl +3

In de populaire cultuur/boeken
Cath Luyten: Het boek Rock Bottom van Cath Luyten behandelt dit onderwerp.
Motivatie: Veel artikelen richten zich op de motivatie om na een dieptepunt weer op te staan.

Borgerhoff \& Lamberigts +2

Kortom, "rock bottom" wordt in de psychologie en filosofie vaak gezien als een pijnlijk, maar transformerend moment dat aanzet tot bezinning, verandering en uiteindelijk groei.

Medium
